\subsection{Skill-guided Plan Composition}
\label{sec:composition}
\label{sec:initial-selection}
\begin{figure}[t]
\centering
\includegraphics[width=\columnwidth]{fig_pipeline.pdf}
\caption{AgentGS planning pipeline. Composition turns retrieved Skills into an initial plan; completion inserts a bridging branch from another Skill to resolve an open input. Without an applicable Skill, planning falls back to the tool-only path.}
\label{fig:pipeline}
\Description{A vertical flowchart. Pills on the left list the stages Retrieve, Select Goals, Expand, Review Support, Search Bridges, and Insert Branch, grouped into Skill-guided Plan Composition and Automatic Plan Completion. Beside each stage a thumbnail shows the plan: retrieved Skills A, B, and D; goal operations of A and B; expanded fragments with an unresolved input; an accepted and a rejected supporting operation; Skill C bridging the fragments; and the completed chain A, C, B. The plan is then grounded over GeoToolKG and executed to produce the answer. A dashed path leads from Retrieve directly to GeoToolKG when no Skill applies.}
\end{figure}

Given a question and input resources, AgentGS retrieves Skills using the requested operations, domain, and available data types. The model selects goal operations from the retrieved candidates. The framework expands their mandatory predecessors and combines the selected fragments into an initial plan, preserving dependencies and permitted tools. Multiple required inputs can produce a branching plan.

This procedure heuristically prunes the planning space. Task relevance narrows the candidate strategies, and goal selection excludes unrelated optional operations while retaining required predecessors. Typed inputs and outputs guide connections between fragments. A raster-producing operation, for example, cannot directly supply a comparison that expects a scalar; the unresolved input identifies a gap for completion.

Each operation retains its source Skill, required inputs, and permitted implementations. These records constrain tool selection and preserve the operation's role in the workflow.

The plan distinguishes data dependencies from control dependencies. Data dependencies identify which earlier result an operation consumes; control dependencies specify ordering without making the preceding output an input. This distinction is retained when fragments from different Skills are composed.

The model thus receives a structured starting plan and scoped tool candidates. It can reuse part of a Skill without regenerating its declared prerequisites or adopting the whole workflow. When no Skill applies, the agent falls back to the existing tool-only planning path.
